\documentclass[12pt]{article}
\usepackage[english,greek]{babel}
\usepackage{amsmath}
\usepackage{enumitem}
\begin{document}
\centering
\title{\selectlanguage{english}Homework 1}
\author{\selectlanguage{greek}Ονοματεπώνυμο: Ελευθέριος Παντζαρτζής \\  ΑΕΜ: 4344}
\date{\today}
\maketitle
\begin{description}
   \item [\selectlanguage{english}(i)] \selectlanguage{greek}
   Η πιθανότητα $Pr[C='\beta']$ είναι: $Pr[C='\beta']=Pr[M='\alpha']*Pr[K=1]+Pr[M='\omega']*Pr[K=2]$ εφόσον μόνο για αυτές τις τιμές των κλειδιών παίρνουμε την τιμή \selectlanguage{english} ciphertext \selectlanguage{greek} β. Και επειδή όλες οι περιπτώσεις κλειδιών είναι ισοπίθανες έχουμε: $Pr[K=1]=Pr[K=2]=1/24$.\\Άρα: $Pr[C='\beta']=0.7*1/24+0.3*1/24=1/24=0.041\overline{6}$
   \item [\selectlanguage{english}(ii)]\selectlanguage{greek}
   Ξανά η πιθανότητα $Pr[C='\theta\pi\delta']$ είναι: $Pr[C='\theta\pi\delta']=Pr[C='\epsilon\nu\alpha']*Pr[K=3]+Pr[M='\delta\upsilon\o']*0$(διότι δεν υπάρχει κανένα κλειδί που να μετατρέπει το $\delta\upsilon\o$ σε $\theta\pi\delta$).\\
   Άρα: $Pr[C='\theta\pi\delta']=0.3*1/24+0=0.0125$.
   \bigbreak
\end{description}
\end{document}